\documentclass[12pt, a4paper]{article} % Clase de documento: 'article' es estándar para trabajos cortos. 12pt es el tamaño de letra. A4 es el formato de papel.

% --- PAQUETES ESENCIALES ---
\usepackage[utf8]{inputenc} % Para la codificación de caracteres, esencial para acentos y eñes
\usepackage[T1]{fontenc}    % Fuente moderna para PDFs
\usepackage[spanish]{babel} % Soporte para idioma español (guionización, fechas, etc.)
\usepackage{amsmath}        % Para matemáticas avanzadas (si fuera necesario)
\usepackage{graphicx}       % Para insertar imágenes
\usepackage{booktabs}       % Mejora el aspecto de las tablas (líneas más profesionales)
\usepackage{geometry}       % Para configurar los márgenes
\usepackage{hyperref}       % Para crear enlaces en el índice y referencias (¡muy recomendado!)
\usepackage{fancyhdr}       % Para personalizar encabezados y pies de página (opcional)
\usepackage[spanish]{babel}
% ¡Asegúrate de que xcolor esté aquí!
\usepackage[table]{xcolor} % 'table' es útil para compatibilidad con listings/tablas
\usepackage{listings}       % Paquete para el formato de código fuente
% --- MODIFICACIÓN DE LA ETIQUETA 'LISTING' A 'LISTADO' ---
\addto{\captionsspanish}{\renewcommand{\lstlistingname}{Listado}}
\addto{\captionsspanish}{\renewcommand{\lstlistlistingname}{Listado de Códigos}}
% --- DEFINICIÓN DE COLOR GRIS Y ESTILO DE CÓDIGO ---

% Definimos un color gris claro para el fondo (usando el modelo RGB)
% Los valores entre 0 y 1. (0.95, 0.95, 0.95 es un gris muy claro)
\definecolor{GrisClaro}{rgb}{0.95, 0.95, 0.95}

% Configuramos el estilo por defecto para todos los bloques de código (lstlisting)
\lstset{
    language=TeX, % Lenguaje por defecto (puedes cambiarlo, ej: Python, C, Java)
    backgroundcolor=\color{GrisClaro}, % Fondo gris claro
    basicstyle=\ttfamily\small,      % Tipo de fuente: typewriter (\ttfamily) y pequeño (\small)
    breaklines=true,                 % Permite saltos de línea largos
    showstringspaces=false,          % No mostrar espacios como caracteres
    numbers=left,                    % Numeración de líneas a la izquierda
    numberstyle=\tiny\color{gray},   % Estilo del número de línea
    frame=single,                    % Dibuja un marco alrededor del código
    framesep=5pt,                    % Espacio entre el marco y el código
    rulesepcolor=\color{black},      % Color de las líneas del marco
    keywordstyle=\color{blue}\bfseries, % Palabras clave en azul y negrita
    commentstyle=\color{green!60!black}\itshape, % Comentarios en verde oscuro y cursiva
    stringstyle=\color{orange},      % Cadenas de texto en naranja
}

% --- LENGUAJES PERSONALIZADOS PARA listings (Dart, Kotlin) ---
\lstdefinelanguage{Dart}{
    keywords={import, class, abstract, static, final, const, void, return, if, else, async, await, Future, bool, String, int, double, true, false, null, this, super, extends, new, try, catch, on, in, is, as, override, enum, switch, case, default, late, factory, required, with, implements, mixin},
    keywordstyle=\color{blue}\bfseries,
    ndkeywords={Widget, State, StatefulWidget, StatelessWidget, BuildContext, MaterialApp, ChangeNotifier, Provider, ChangeNotifierProvider, Consumer, MultiProvider, Scaffold, AppBar, Column, Row, Container, Text, Icon, GestureDetector, ApiClient, AuthRepository, LoginViewModel},
    ndkeywordstyle=\color{purple}\bfseries,
    sensitive=true,
    comment=[l]{//},
    morecomment=[s]{/*}{*/},
    morestring=[b]",
    morestring=[b]',
}
\lstdefinelanguage{Kotlin}{
    keywords={package, import, fun, override, val, var, when, class, return, super, private, public, internal, object, this, if, else, true, false, null, in, is, as},
    keywordstyle=\color{blue}\bfseries,
    ndkeywords={String, Int, Boolean, Unit, Bundle, MethodChannel, FlutterFragmentActivity, FlutterEngine, WindowManager},
    ndkeywordstyle=\color{purple}\bfseries,
    sensitive=true,
    comment=[l]{//},
    morecomment=[s]{/*}{*/},
    morestring=[b]",
    morestring=[b]',
}

% --- CONFIGURACIÓN DE MARGENES ---
\geometry{
 a4paper,
 total={170mm,257mm},
 left=25mm,
 right=25mm,
 top=30mm,
 bottom=30mm,
 }

% --- CONFIGURACIÓN DE ENCABEZADOS Y PIES (OPCIONAL) ---
\pagestyle{fancy}
\fancyhf{} % Limpiar encabezados y pies
\renewcommand{\headrulewidth}{0pt} % Eliminar línea en el encabezado
\rfoot{\thepage} % Número de página a la derecha en el pie
\lfoot{ Desarrollo de Aplicaciones Móviles con Flutter / Miguel Ángel Gutiérrez Gómez} % Texto en el pie de página izquierdo

% --- DATOS DEL DOCUMENTO (PARA LA PORTADA) ---
\newcommand{\nombreuniversidad}{Universidad Politécnica de Chiapas}
\newcommand{\nombrecarrera}{Ingeniería en Software}
\newcommand{\nombremateria}{Desarrollo de Aplicaciones Móviles con Flutter}
\newcommand{\nombreprofesor}{MC. José Alonso Macías Montoya}
\newcommand{\corte}{Corte 1}
\newcommand{\actividad}{Aplicación Transaccional - VisionPrice}
\newcommand{\nombres}{Miguel Ángel Gutiérrez Gómez}
\newcommand{\matricula}{233382}
\newcommand{\fechaentrega}{\today} % \today inserta la fecha actual automáticamente
\newcommand{\ciudad}{Suchiapa, Chiapas}

\usepackage[utf8]{inputenc} % ¡ESTA LÍNEA ES FUNDAMENTAL!
\usepackage[T1]{fontenc}
\usepackage[spanish]{babel}
% --- INICIO DEL DOCUMENTO ---
\begin{document}

% --- PORTADA ---
\begin{titlepage}
    \centering
    \vspace*{1cm} % Espacio vertical
    {\Huge\bfseries \nombreuniversidad \par} % Título principal: Nombre de la Universidad
    \vspace{0.5cm}
    {\Large \nombrecarrera \par} % Nombre de la Carrera

    \vspace{2cm}
    \rule{\textwidth}{1.5pt} % Línea gruesa
    \vspace{0.4cm}
    {\Large\bfseries \actividad \par} % Título del Trabajo
    \vspace{0.4cm}
    \rule{\textwidth}{1.5pt} % Línea gruesa

    \vspace{3cm}

    % --- DATOS DEL ALUMNO Y CURSO (EN UNA SOLA COLUMNA) ---
    % Se eliminaron los entornos minipage y \hfill
    \begin{flushleft}
        \centering % Centra las líneas dentro del entorno flushleft
        \textbf{Materia:} \nombremateria \\
        \textbf{Profesor:} \nombreprofesor \\
        \textbf{Corte:} \corte
    \end{flushleft}

    \vspace{0.5cm} % Pequeño espacio de separación

    \begin{flushleft}
        \centering % Centra las líneas dentro del entorno flushleft
        \textbf{Alumno:} \nombres \\
        \textbf{Matrícula:} \matricula \\
    \end{flushleft}
    % --- FIN DE DATOS EN UNA SOLA COLUMNA ---

    \vfill % Empuja el contenido restante al final de la página

    % Fecha y Ciudad
    \vspace{1cm}
    \textbf{\ciudad, a \fechaentrega}

\end{titlepage}

\clearpage % Salto de página para empezar el contenido

% --- ÍNDICE/TABLA DE CONTENIDO ---
% Se recomienda usar la instrucción \tableofcontents inmediatamente después de la portada.
% LaTeX genera automáticamente el índice a partir de las secciones (\section, \subsection, etc.)
\tableofcontents
\clearpage

% --- CUERPO DEL TRABAJO ---

\section{Introducción}
\label{sec:introduccion}
VisionPrice es una aplicación móvil \textbf{transaccional} construida con
Flutter cuyo objetivo es gestionar de forma profesional los presupuestos
de proyectos de construcción y remodelación. La app permite a un
contratista (o vendedor de materiales) iniciar sesión de forma segura,
crear nuevos proyectos asignándoles cliente, ciudad y tipo de trabajo
(piso, pared, techo o combinado), consultar, editar y eliminar los ya
existentes, y revisar las métricas agregadas (proyectos activos, en
cotización) desde un dashboard.

El propósito de este trabajo es documentar la arquitectura y los
patrones aplicados en la construcción de la aplicación, demostrando que
cumple con los requisitos académicos planteados en la rúbrica
\textit{App Transaccional}: \textbf{Arquitectura} (MVVM + Clean +
Screaming + Vertical Slicing), \textbf{Gestión de estado} con Provider,
\textbf{Cliente HTTP} con verbos GET/POST/PUT/DELETE, \textbf{Interfaz
gráfica} basada en Material Theme 3.0 y Atomic Design, y consumo de una
\textbf{API RESTful} dividida en dos microservicios (auth y projects).

\section{Recursos Utilizados}
\label{sec:recursos}
En esta sección se listan las herramientas, librerías y servicios
fundamentales para la realización de la actividad.

\begin{itemize}
    \item \textbf{Lenguaje y SDK:} Dart 3.12 y Flutter 3.44 para el
    cliente móvil; Node.js 20 + TypeScript 5 para el backend.
    \item \textbf{Librerías Flutter:} \texttt{provider} (gestión de
    estado), \texttt{http} (cliente REST),
    \texttt{flutter\_secure\_storage} (Keystore/Keychain),
    \texttt{google\_sign\_in} 7.x (login federado),
    \texttt{local\_auth} 3.x (huella / Face ID),
    \texttt{geolocator} (detección de Fake GPS).
    \item \textbf{Backend:} Fastify 4, Prisma 5, MySQL 8 en
    contenedores Docker, JWT (\texttt{jsonwebtoken}), bcrypt, Zod.
    \item \textbf{Herramientas:} Android Studio, VS Code, Postman,
    Docker Desktop, Git, Overleaf y compilador \LaTeX{}.
    \item \textbf{Diseño:} Material Theme 3.0, prototipo VisionPrice
    (HTML standalone) y guía de estilo Atomic Design.
\end{itemize}

\section{Desarrollo}
\label{sec:desarrollo}
Esta sección describe la metodología, la arquitectura, el manejo de
estado, el cliente HTTP y la interfaz gráfica de la aplicación.

\subsection{Arquitectura}
\label{sec:arquitectura}
La aplicación combina cuatro principios complementarios.

\subsubsection{Clean Architecture}
La capa de dominio (entidades, value objects, repositorios abstractos y
casos de uso) es \textbf{independiente} de cualquier framework. Toda
dependencia apunta hacia adentro: \texttt{presentation} depende de
\texttt{domain}, \texttt{data} depende de \texttt{domain}, y
\texttt{domain} no depende de nadie.

\subsubsection{Screaming Architecture}
Al abrir el árbol del proyecto, la estructura \textit{grita} lo que hace
la aplicación, no la tecnología empleada. Las carpetas de primer nivel
son \texttt{auth} y \texttt{projects}, no \texttt{controllers} ni
\texttt{models}.

\subsubsection{Vertical Slicing}
Cada feature contiene \textbf{toda} su pila: \texttt{data},
\texttt{domain}, \texttt{presentation} y \texttt{di}. Esto desacopla
features entre sí y permite trabajar en uno sin afectar al resto.

\subsubsection{MVVM}
Cada pantalla tiene un \texttt{ViewModel} (\texttt{ChangeNotifier}) con
un \texttt{State} inmutable. La \texttt{View} solo construye widgets y
reacciona a cambios; las acciones del usuario llaman métodos del
ViewModel, que invoca los casos de uso del dominio.

\begin{lstlisting}[language={}, caption={Estructura de carpetas (Screaming + Vertical Slicing)}]
lib/
  core/                  (UI compartida + utilidades transversales)
  features/
    auth/                (slice vertical: autenticacion)
      data/
        datasources/{remote,local}/
        local/
        mappers/
        repositories/
      domain/
        entities/
        value_objects/
        services/        (ports: BiometricAuthenticator, CredentialsVault)
        repositories/    (port: AuthRepository)
        use_cases/
      presentation/
        pages/
        view_models/
      di/auth_module.dart
    projects/            (slice vertical: proyectos)
      data/ domain/ presentation/ di/
\end{lstlisting}

\begin{lstlisting}[language=Dart, caption={ViewModel MVVM con estado inmutable}]
class LoginViewModel extends ChangeNotifier {
  final LoginWithPassword _loginWithPassword;
  LoginState _state = const LoginState();
  LoginState get state => _state;

  Future<bool> loginWithPassword({
    required String email,
    required String password,
  }) async {
    if (_state.isBusy) return false;
    _set(_state.copyWith(isSubmitting: true, errorMessage: null));
    try {
      await _loginWithPassword(email: email, password: password);
      return true;
    } on ApiException catch (e) {
      _set(_state.copyWith(errorMessage: e.message));
      return false;
    } finally {
      _set(_state.copyWith(isSubmitting: false));
    }
  }

  void _set(LoginState next) { _state = next; notifyListeners(); }
}
\end{lstlisting}

\subsection{Gestión de Estado con Provider}
\label{sec:provider}
La aplicación utiliza la librería \texttt{provider} con dos niveles de
provisión:

\begin{itemize}
    \item \textbf{Nivel global:} los módulos
    (\texttt{AuthModule}, \texttt{ProjectsModule}) son singletons
    inyectados con \texttt{Provider.value} en el composition root. Las
    páginas los leen con \texttt{context.read<AuthModule>()}.
    \item \textbf{Nivel por página:} cada \texttt{Page} envuelve su
    árbol en un \texttt{ChangeNotifierProvider} que construye su
    ViewModel con la factoría del módulo, garantizando una instancia
    fresca por pantalla y evitando fugas de estado entre
    navegaciones.
\end{itemize}

\begin{lstlisting}[language=Dart, caption={Composition root global con MultiProvider}]
return MultiProvider(
  providers: [
    Provider<AuthModule>.value(value: _authModule),
    Provider<ProjectsModule>.value(value: _projectsModule),
  ],
  child: MaterialApp(
    title: 'VisionPrice',
    theme: AppTheme.light(),
    initialRoute: Routes.integrityGate,
    routes: {
      Routes.login:         (_) => const LoginPage(),
      Routes.register:      (_) => const RegisterPage(),
      Routes.projects:      (_) => const DashboardPage(),
      Routes.addProject:    (_) => const AddProjectPage(),
      Routes.editProject:   (_) => const EditProjectPage(),
      Routes.projectDetail: (_) => const ProjectDetailPage(),
    },
  ),
);
\end{lstlisting}

\begin{lstlisting}[language=Dart, caption={ChangeNotifierProvider + Consumer en la página}]
class LoginPage extends StatelessWidget {
  @override
  Widget build(BuildContext context) {
    final module = context.read<AuthModule>();
    return ChangeNotifierProvider<LoginViewModel>(
      create: (_) => module.loginViewModelFactory(),
      child: const _LoginView(),
    );
  }
}

return Consumer<LoginViewModel>(
  builder: (context, vm, _) {
    final state = vm.state;
    return PrimaryButton(
      label: 'Continuar',
      loading: state.isSubmitting,
      onPressed: state.isBusy ? null : () => _onContinue(vm),
    );
  },
);
\end{lstlisting}

\subsection{Cliente HTTP}
\label{sec:http}
La aplicación utiliza el paquete oficial \texttt{http}. Internamente se
construyó un wrapper \texttt{ApiClient} que centraliza la URL base, los
headers JSON, el header \texttt{Authorization: Bearer <token>}, el
timeout y la decodificación de errores. A continuación se muestran
ejemplos reales de los cuatro verbos REST consumidos por la app.

\begin{lstlisting}[language=Dart, caption={GET /projects: listar proyectos}]
Future<List<ProjectDto>> list({String? status}) async {
  final token = await _requireToken();
  final query = status == null ? '' : '?status=$status';
  final json = await _apiClient.getJson(
    '/projects$query',
    accessToken: token,
  );
  final raw = (json['projects'] as List<dynamic>?) ?? const [];
  return raw
      .map((e) => ProjectDto.fromJson(e as Map<String, dynamic>))
      .toList(growable: false);
}
\end{lstlisting}

\begin{lstlisting}[language=Dart, caption={POST /auth/login: iniciar sesión}]
Future<AuthResponseDto> loginWithPassword({
  required String email,
  required String password,
}) async {
  final json = await _apiClient.postJson('/auth/login', {
    'email': email,
    'password': password,
  });
  return AuthResponseDto.fromJson(json);
}
\end{lstlisting}

\begin{lstlisting}[language=Dart, caption={PUT/PATCH /projects/:id: actualizar proyecto}]
Future<ProjectDto> update(String id, Map<String, dynamic> changes) async {
  final token = await _requireToken();
  final res = await _apiClient.patch('/projects/$id', changes,
      accessToken: token);
  if (res.statusCode < 200 || res.statusCode >= 300) {
    throw _apiClient.decodeError(res);
  }
  final body = _apiClient.decodeJson(res);
  return ProjectDto.fromJson(body['project'] as Map<String, dynamic>);
}
\end{lstlisting}

\begin{lstlisting}[language=Dart, caption={DELETE /projects/:id: eliminar proyecto}]
Future<void> delete(String id) async {
  final token = await _requireToken();
  final res = await _apiClient.deleteRaw('/projects/$id',
      accessToken: token);
  if (res.statusCode == 204 || res.statusCode == 200) return;
  throw _apiClient.decodeError(res);
}
\end{lstlisting}

\subsection{Interfaz Gráfica}
\label{sec:ui}
Toda la app comparte un \texttt{AppTheme} con \texttt{useMaterial3: true},
paleta consistente, tipografía, botones, inputs y checkboxes
unificados. La capa \texttt{core/ui} sigue \textbf{Atomic Design}:
\textit{atoms} (botones, etiquetas, iconos), \textit{molecules} (campos
con label, password con toggle, city picker) y \textit{organisms}
(brand header, social auth buttons).

Los widgets más representativos usados a lo largo de la app incluyen:
\texttt{Scaffold}, \texttt{AppBar}, \texttt{SafeArea},
\texttt{SingleChildScrollView}, \texttt{Column}, \texttt{Row},
\texttt{Form}, \texttt{TextFormField}, \texttt{GridView.builder},
\texttt{ChoiceChip}, \texttt{FloatingActionButton},
\texttt{AlertDialog}, \texttt{SnackBar}, \texttt{Consumer},
\texttt{ChangeNotifierProvider}, además de los átomos y moléculas
propios del proyecto.

\begin{figure}[h]
    \centering
    \fbox{\parbox[c][6cm][c]{0.55\textwidth}{\centering \textit{Pantalla de Login (email/password, Google y Face ID)}}}
    \caption{Pantalla de Login}
    \label{fig:login}
\end{figure}

\begin{figure}[h]
    \centering
    \fbox{\parbox[c][6cm][c]{0.55\textwidth}{\centering \textit{Pantalla de Registro}}}
    \caption{Pantalla de Registro}
    \label{fig:register}
\end{figure}

\begin{figure}[h]
    \centering
    \fbox{\parbox[c][6cm][c]{0.55\textwidth}{\centering \textit{Dashboard con KPIs y lista de proyectos}}}
    \caption{Dashboard principal}
    \label{fig:dashboard}
\end{figure}

\begin{figure}[h]
    \centering
    \fbox{\parbox[c][6cm][c]{0.55\textwidth}{\centering \textit{Nuevo proyecto: pasos numerados + grid 2x2}}}
    \caption{Formulario de creación de proyecto}
    \label{fig:add-project}
\end{figure}

\begin{figure}[h]
    \centering
    \fbox{\parbox[c][6cm][c]{0.55\textwidth}{\centering \textit{Detalle del proyecto con editar/eliminar}}}
    \caption{Detalle de proyecto}
    \label{fig:detail}
\end{figure}

\subsection{Endpoints REST consumidos}
La aplicación consume los endpoints REST resumidos en la
Tabla \ref{tab:endpoints}.

\begin{table}[h]
    \centering
    \caption{Endpoints REST consumidos por la aplicación}
    \label{tab:endpoints}
    \begin{tabular}{llll}
        \toprule
        \textbf{Método} & \textbf{Ruta} & \textbf{Auth} & \textbf{Propósito} \\
        \midrule
        POST   & /auth/register     & público & Crear cuenta \\
        POST   & /auth/login        & público & Iniciar sesión \\
        POST   & /auth/google/token & público & Login Google idToken \\
        POST   & /auth/refresh      & público & Rotar tokens \\
        POST   & /auth/logout       & Bearer  & Revocar sesión \\
        GET    & /auth/me           & Bearer  & Perfil actual \\
        GET    & /projects          & Bearer  & Listar proyectos \\
        POST   & /projects          & Bearer  & Crear proyecto \\
        GET    & /projects/:id      & Bearer  & Detalle proyecto \\
        PATCH  & /projects/:id      & Bearer  & Actualizar proyecto \\
        DELETE & /projects/:id      & Bearer  & Eliminar proyecto \\
        \bottomrule
    \end{tabular}
    \vspace{0.5em}
    \footnotesize Nota: PATCH se utiliza en lugar de PUT para permitir actualizaciones parciales del recurso.
\end{table}

\subsection{Prompts Utilizados}
\label{sec:prompts}
A continuación se presentan los cinco prompts más representativos
utilizados con un asistente LLM (Claude) durante el desarrollo. Cada
respuesta del asistente fue revisada, modificada e integrada con
criterio propio; el asistente actuó como \textit{pair programmer},
nunca como autor único.

\subsubsection{Prompt 1 — Backend con arquitectura hexagonal}
\textit{``Scaffold a production-ready authentication microservice using
Node.js + TypeScript + Fastify following strict Hexagonal Architecture
(Ports \& Adapters) combined with Clean Architecture: domain (entities
without ORM, value objects, repository ports, errors), application
(use cases, DTOs), infrastructure (Prisma, JWT, Google), interfaces
(HTTP). Manual DI in main.ts. Env vars validated with Zod''.}

Resultado: \texttt{auth-service} completo con 6 endpoints y arquitectura
hexagonal estricta.

\subsubsection{Prompt 2 — Frontend con Clean + MVVM + Atomic Design}
\textit{``Adaptar lo que tengo de Flutter al estilo Clean Architecture +
MVVM + Repository + Atomic Design, con todas las mejoras recomendadas:
atomic design completo, sin código duplicado entre features, value
objects en domain, repository real con use cases, sin fugas de
dependencia. State management con ChangeNotifier + Provider. DI
manual''.}

Resultado: reorganización completa del frontend en \texttt{core/} y
\texttt{features/}, atoms y molecules compartidos, ViewModels MVVM
puros.

\subsubsection{Prompt 3 — CRUD completo con microservicio separado}
\textit{``Construye el feature de Projects: tanto en el frontend como en
el backend. Microservicio nuevo VisionPriceProjects con cliente +
ubicación + presupuesto + workType + area + status. CRUD completo: Lista
+ Crear + Detalle + Editar + Eliminar. Mismo JWT secret entre los dos
servicios''.}

Resultado: \texttt{projects-service} en puerto 3001 que verifica los
JWT emitidos por \texttt{auth-service}, y feature \texttt{projects} en
Flutter con todas las pantallas.

\subsubsection{Prompt 4 — Adaptación al diseño del prototipo}
\textit{``Adapta lo que está al diseño del prototipo VisionPrice:
Dashboard con header (fecha + saludo + avatar), KPIs ACTIVOS y EN
COTIZACIÓN, lista de proyectos con thumbnail, ciudad, tiempo relativo y
status chip; formulario Nuevo proyecto con 3 pasos numerados + grid 2x2
de tipos de trabajo + city picker + banner de precios actualizados''.}

Resultado: rediseño completo del Dashboard, AddProjectPage con pasos
numerados, nuevos átomos (MetricCard, NumberedStepLabel,
WorkTypeChoiceCard, CityPicker, PricesUpdateBanner).

\subsubsection{Prompt 5 — Acceso rápido por biometría}
\textit{``Implementa autenticación biométrica (Face ID / huella) siguiendo
Clean Architecture. Las credenciales deben quedar guardadas de forma
segura cuando el usuario active el acceso rápido, y solo deben cambiar
cuando un usuario distinto inicie sesión y vuelva a activarlo. Para
Google, evita que aparezca el picker después de la biometría''.}

Resultado: dos nuevos ports en \texttt{domain}
(\texttt{BiometricAuthenticator}, \texttt{CredentialsVault}), adapters
con \texttt{flutter\_secure\_storage}, casos de uso y cascada
\textit{refresh $\to$ silent Google $\to$ picker} en
\texttt{loginWithBiometric}.

\section{Resultados y Conclusiones}
\label{sec:resultados}
En esta sección se presentan los \textbf{resultados clave} y las
\textbf{conclusiones} obtenidas tras el desarrollo.

\begin{itemize}
    \item \textbf{Resultado 1:} Se entregó una aplicación móvil
    transaccional funcional con CRUD completo de proyectos sobre dos
    microservicios REST y autenticación con email/password, Google y
    biometría.
    \item \textbf{Resultado 2:} Se aplicó de forma estricta
    \textbf{Clean Architecture}, \textbf{Screaming Architecture},
    \textbf{Vertical Slicing} y \textbf{MVVM}, manteniendo las capas
    desacopladas (la capa de dominio no importa ningún framework).
    \item \textbf{Resultado 3:} Se utilizó \textbf{Provider} en dos
    niveles (módulos globales + ViewModel por página) y un \texttt{State}
    inmutable con \texttt{copyWith}, eliminando bugs de mutación
    accidental.
    \item \textbf{Resultado 4:} Se construyó un \texttt{ApiClient}
    reutilizable con los cuatro verbos REST (GET, POST, PATCH, DELETE)
    y manejo centralizado de errores.
    \item \textbf{Conclusión:} El costo en disciplina (más archivos,
    más interfaces) que impone Clean Architecture se recupera con
    creces al refactorizar: cambiar \texttt{google\_sign\_in} v6 a v7 o
    reemplazar widgets que causaban bugs en Flutter 3.44 se redujo a
    tocar \textbf{un solo archivo} en la capa \texttt{data} o
    \texttt{presentation}, sin tocar el dominio.
    \item \textbf{Aprendizaje:} Integrar dos microservicios con JWT
    compartido, biometría con cascada inteligente y MVVM puro en
    Flutter exige planificar las interfaces (ports) antes que los
    frameworks. Cuando esa regla se respeta, el código es testeable,
    mantenible y resiste cambios en las dependencias externas.
\end{itemize}

\clearpage % Salto de página para empezar los anexos

\section{Anexos}
\label{sec:anexos}
Los anexos incluyen material complementario: el repositorio del
proyecto, capturas adicionales y fragmentos de código relevantes.

\subsection{Repositorio del proyecto en GitHub}
\label{sec:repo}
El código fuente completo de la aplicación (cliente Flutter y los dos
microservicios backend) se encuentra publicado en el siguiente
repositorio de GitHub:

\begin{itemize}
    \item \textbf{Repositorio:} \url{https://github.com/MiguelAngelGutierrezGomez/VisionPrice}
    \item \textbf{Rama principal:} \texttt{main}
    \item \textbf{Estructura del workspace:}
\end{itemize}

\begin{lstlisting}[language={}, caption={Layout del workspace en el repositorio}]
VisionPrice/
  VisionPriceAuth/      (auth-service: Node + TS + Fastify + Prisma)
  VisionPriceProjects/  (projects-service)
  VisionPrice/          (cliente Flutter)
  reportes/             (este documento y sus imagenes)
\end{lstlisting}

\subsection{Capturas de pantalla adicionales}
A continuación se muestra una captura adicional del flujo de
autenticación con biometría (huella digital / Face ID) sobre el botón
\textit{Acceso rápido} de la pantalla de Login.

\begin{figure}[h]
    \centering
    \fbox{\parbox[c][6cm][c]{0.55\textwidth}{\centering \textit{Acceso rápido por biometría tras enrolamiento}}}
    \caption{Captura del flujo de acceso rápido por biometría.}
    \label{fig:biometric}
\end{figure}

\clearpage

\subsection{Código fuente adicional}
Para incluir código fuente de manera profesional y con estilo, se
utiliza el entorno \texttt{lstlisting}. A continuación, dos fragmentos
representativos: el \textit{composition root} del módulo de
autenticación (DI manual) y la activación de \texttt{FLAG\_SECURE} en
Android desde Kotlin.

\begin{lstlisting}[language=Dart, caption={Composition root manual del feature Auth}]
class AuthModule {
  final AuthRepository repository;
  final LoginWithPassword loginWithPasswordUseCase;
  final LoginWithGoogle loginWithGoogleUseCase;
  final LoginWithBiometric loginWithBiometricUseCase;

  factory AuthModule.create({required ApiClient apiClient}) {
    final remote = AuthRemoteDataSource(apiClient);
    final storage = TokenStorage();
    final google = GoogleSignInService();
    final BiometricAuthenticator biometric = LocalAuthBiometricAdapter();
    final CredentialsVault vault = CredentialsVaultImpl();
    final AuthRepository repo = AuthRepositoryImpl(
      remote: remote, storage: storage, google: google,
      biometric: biometric, vault: vault,
    );
    return AuthModule._(
      repository: repo,
      loginWithPasswordUseCase: LoginWithPassword(repo),
      loginWithGoogleUseCase: LoginWithGoogle(repo),
      loginWithBiometricUseCase: LoginWithBiometric(repo),
    );
  }

  LoginViewModel loginViewModelFactory() => LoginViewModel(
    loginWithPassword: loginWithPasswordUseCase,
    loginWithGoogle: loginWithGoogleUseCase,
    loginWithBiometric: loginWithBiometricUseCase,
  );
}
\end{lstlisting}

\vspace{1cm}

\begin{lstlisting}[language=Kotlin, caption={Activación de FLAG\_SECURE en MainActivity}]
class MainActivity : FlutterFragmentActivity() {
    private val CHANNEL = "visionprice/secure_screen"

    override fun configureFlutterEngine(flutterEngine: FlutterEngine) {
        super.configureFlutterEngine(flutterEngine)
        MethodChannel(flutterEngine.dartExecutor.binaryMessenger, CHANNEL)
            .setMethodCallHandler { call, result ->
                when (call.method) {
                    "enable" -> {
                        window.setFlags(
                            WindowManager.LayoutParams.FLAG_SECURE,
                            WindowManager.LayoutParams.FLAG_SECURE
                        )
                        result.success(null)
                    }
                    "disable" -> {
                        window.clearFlags(WindowManager.LayoutParams.FLAG_SECURE)
                        result.success(null)
                    }
                    else -> result.notImplemented()
                }
            }
    }
}
\end{lstlisting}

\end{document}
